\section{Chapter 24}

\begin{quotex}
\begin{verse}
The effort of one who stands on tiptoe does not lead to rising\\
To splay the legs excessively is not walking\\
He who puts himself in the light remains in the shadow\\
He who believes he has arrived is pushed back\\
To put oneself ahead is to lose value\\
To esteem oneself is to decay\\
With respect to the Way, all this\\
Is nothing but excrescence and refuse\\
With respect to the goal, it is something to be disdained\\
Far removed from all this is he who is in the Way.
\end{verse}
\end{quotex}

\paragraph{Commentary}


The path of one united with the Principle is marked by natural, impersonal spontaneity, opposed to all behaviors based on shifting and externalizing the center of oneself into the concretion of the small ego, mere excrescence (the text adds: ``refuse'') with respect to being. Any accentuation of the latter automatically converts into a diminishment of the former: a Taoist theme already familiar to us.

In the same line of thought, a still more drastic Taoist expression is: the ego is an abscess that bursts at death.

\paragraph{Chinese text and literal translation}


Chapter 24 (\begin{CJK*}{UTF8}{bsmi}第二十四章\end{CJK*})
\columnratio{0.4}
\begin{paracol}{2}
\begin{verse}
\begin{CJK*}{UTF8}{bsmi}
跂者不立，\\
跨者不行，\\
自見者不明，\\
自是者不彰。\\
自伐者無功，\\
自矜者不長。\\
其在道也，\\
曰餘食贅行。\\
物或惡之，\\
故有道者不處也。
\end{CJK*}
\end{verse}
\switchcolumn
\begin{verse}
Ones who tip-toe do not stand,\\
Ones who stride do not pace,\\
Ones show off themselves do not shine,\\
Ones that bias themselves with self-righteousness do not justify.\\
Ones self-asserted are without accomplishments,\\
Ones self-esteemed are not endurable.\\
These according to the Dao,\\
Are called extravagance and exorbitance.\\
which the masses resent,\\
Therefore those embracing the Dao dwell not.
\end{verse}
\end{paracol}

\flrightit{Posted on 2025-06-04 by Tannheuser}

\begin{center}* * *\end{center}

\begin{footnotesize}
\begin{sffamily}
\texttt{Luka Markovic on 2025-10-18 at 14:56 said:}

Chapter 24 contrasts The Way, or the process of subtle inner maturation, with being merely intellectually acquisitive. The germ of wisdom ripens through cultivation.

\end{sffamily}
\end{footnotesize}

